
\documentclass{article}
\usepackage{colortbl}
\usepackage{makecell}
\usepackage{multirow}
\usepackage{supertabular}

\begin{document}

\newcounter{utterance}

\centering \large Interaction Transcript for game `function\_detective', experiment `two\_numbers\_passive\_examples', episode 289 with gpt{-}5.4{-}mini.
\vspace{24pt}

{ \footnotesize  \setcounter{utterance}{1}
\setlength{\tabcolsep}{0pt}
\begin{supertabular}{c@{$\;$}|p{.15\linewidth}@{}p{.15\linewidth}p{.15\linewidth}p{.15\linewidth}p{.15\linewidth}p{.15\linewidth}}
    \# & \multicolumn{2}{c}{Player} && \multicolumn{2}{c}{Game Master} \\
    \hline

    \theutterance \stepcounter{utterance}  
    & & & \multicolumn{4}{p{0.6\linewidth}}{
        \cellcolor[rgb]{0.9,0.9,0.9}{
            \makecell[{{p{\linewidth}}}]{
                \texttt{\tiny{[P1$\langle$GM]}}
                \texttt{I have a hidden Python function for you to figure out.} \\
\\ 
\texttt{Your job is to infer what it does and then submit a correct Python implementation.} \\
\\ 
\texttt{Function signature:} \\
\texttt{(a: int, b: int) {-}> int} \\
\\ 
\texttt{You have 15 turns total.} \\
\\ 
\texttt{I control which valid input{-}output examples you see. You can request them one at a time.} \\
\\ 
\texttt{You must make your best SOLVE attempt no later than the second{-}to{-}last turn.} \\
\texttt{Do not wait until the very end.} \\
\\ 
\texttt{You can use these actions:} \\
\texttt{{-} JSON format described at the end of the prompt} \\
\texttt{{-} SOLVE: ```python ... ```} \\
\\ 
\texttt{Do not add import statements in your guess.} \\
\texttt{Anything you might need is already available, so just write the function body in hypothesis and return the answer.} \\
\\ 
\texttt{Preloaded examples:} \\
\texttt{1. (({-}99, 10), 723)} \\
\\ 
\texttt{Scoring:} \\
\texttt{{-} correctness is required for any non{-}zero score} \\
\texttt{{-} in interactive modes, solving with fewer information requests is better} \\
\\ 
\texttt{Input formatting reminders:} \\
\texttt{{-} strings must be quoted} \\
\texttt{{-} booleans must be True/False} \\
\texttt{{-} use commas between multiple values} \\
\texttt{{-} do not write any extra action tags in the same message} \\
\\ 
\texttt{Format example:} \\
\texttt{You write:} \\
\texttt{\{'current\_mode': <probing/confirming/don't know; choose whichever best describes your current state in the task>, 'input\_rationale': <brief description of why you chose this hypothesis>, 'hypothesis': <either the exact string "unknown" if you have no guess yet OR your current best guess as runnable Python matching the given signature, with a one{-}line docstring stating what it does. Never write placeholder code, if you cannot commit to a guess, use "unknown". Match the given signature.} \\
\texttt{e.g.} \\
\texttt{def f(x):} \\
\texttt{	"""Returns the absolute value of x."""} \\
\texttt{	return x if x >= 0 else {-}x} \\
\\ 
\texttt{Use "unknown" if you have no guess.>, 'confidence\_score': <score between 1{-}100. 100 means highest confidence in your hypothesis>\}} \\
\\ 
\texttt{Output only one valid json per turn.} \\
            }
        }
    }
     \\ \\

    \theutterance \stepcounter{utterance}  
    & \multicolumn{4}{p{0.6\linewidth}}{
        \cellcolor[rgb]{0.561,0.737,0.561}{
            \makecell[{{p{\linewidth}}}]{
                \texttt{\tiny{[P1$\rangle$GM]}}
                \texttt{\{"current\_mode":"probing","input\_rationale":"I only have one example, which is insufficient to identify the rule. I need a few more input{-}output pairs to infer the hidden function reliably.","hypothesis":"unknown","confidence\_score":5\}} \\
            }
        }
    }
    & & \\ \\

    \theutterance \stepcounter{utterance}  
    & & & \multicolumn{4}{p{0.6\linewidth}}{
        \cellcolor[rgb]{0.9,0.9,0.9}{
            \makecell[{{p{\linewidth}}}]{
                \texttt{\tiny{[P1$\langle$GM]}}
                \texttt{OUTPUT: (({-}1, {-}99), {-}399)} \\
            }
        }
    }
     \\ \\

    \theutterance \stepcounter{utterance}  
    & \multicolumn{4}{p{0.6\linewidth}}{
        \cellcolor[rgb]{0.561,0.737,0.561}{
            \makecell[{{p{\linewidth}}}]{
                \texttt{\tiny{[P1$\rangle$GM]}}
                \texttt{\{"current\_mode":"probing","input\_rationale":"With only two examples, there are many possible integer formulas. I need at least one more pair to distinguish among linear and non{-}linear candidates.","hypothesis":"unknown","confidence\_score":8\}} \\
            }
        }
    }
    & & \\ \\

    \theutterance \stepcounter{utterance}  
    & & & \multicolumn{4}{p{0.6\linewidth}}{
        \cellcolor[rgb]{0.9,0.9,0.9}{
            \makecell[{{p{\linewidth}}}]{
                \texttt{\tiny{[P1$\langle$GM]}}
                \texttt{OUTPUT: (({-}99, {-}2), {-}501)} \\
            }
        }
    }
     \\ \\

    \theutterance \stepcounter{utterance}  
    & \multicolumn{4}{p{0.6\linewidth}}{
        \cellcolor[rgb]{0.561,0.737,0.561}{
            \makecell[{{p{\linewidth}}}]{
                \texttt{\tiny{[P1$\rangle$GM]}}
                \texttt{\{"current\_mode":"confirming","input\_rationale":"The three examples fit a simple affine function exactly. Solving the linear system gives f(a,b)=102a+102b+9801, which matches all observed outputs.","hypothesis":"def f(a: int, b: int) {-}> int:\textbackslash{}n\textbackslash{}t\textbackslash{}"\textbackslash{}"\textbackslash{}"Returns 102 times the sum of the inputs plus 9801.\textbackslash{}"\textbackslash{}"\textbackslash{}"\textbackslash{}n\textbackslash{}treturn 102 * (a + b) + 9801","confidence\_score":97\}} \\
            }
        }
    }
    & & \\ \\

    \theutterance \stepcounter{utterance}  
    & & & \multicolumn{4}{p{0.6\linewidth}}{
        \cellcolor[rgb]{0.9,0.9,0.9}{
            \makecell[{{p{\linewidth}}}]{
                \texttt{\tiny{[P1$\langle$GM]}}
                \texttt{OUTPUT: (({-}99, 2), {-}93)} \\
            }
        }
    }
     \\ \\

    \theutterance \stepcounter{utterance}  
    & \multicolumn{4}{p{0.6\linewidth}}{
        \cellcolor[rgb]{0.561,0.737,0.561}{
            \makecell[{{p{\linewidth}}}]{
                \texttt{\tiny{[P1$\rangle$GM]}}
                \texttt{\{"current\_mode":"confirming","input\_rationale":"The new example also matches the same affine rule, strengthening the hypothesis that the function is 102*(a+b)+9801.","hypothesis":"def f(a: int, b: int) {-}> int:\textbackslash{}n\textbackslash{}t\textbackslash{}"\textbackslash{}"\textbackslash{}"Returns 102 times the sum of the inputs plus 9801.\textbackslash{}"\textbackslash{}"\textbackslash{}"\textbackslash{}n\textbackslash{}treturn 102 * (a + b) + 9801","confidence\_score":99\}} \\
            }
        }
    }
    & & \\ \\

    \theutterance \stepcounter{utterance}  
    & & & \multicolumn{4}{p{0.6\linewidth}}{
        \cellcolor[rgb]{0.9,0.9,0.9}{
            \makecell[{{p{\linewidth}}}]{
                \texttt{\tiny{[P1$\langle$GM]}}
                \texttt{OUTPUT: (({-}10, {-}1), {-}43)} \\
            }
        }
    }
     \\ \\

    \theutterance \stepcounter{utterance}  
    & \multicolumn{4}{p{0.6\linewidth}}{
        \cellcolor[rgb]{0.561,0.737,0.561}{
            \makecell[{{p{\linewidth}}}]{
                \texttt{\tiny{[P1$\rangle$GM]}}
                \texttt{\{"current\_mode":"probing","input\_rationale":"The new example breaks the current affine hypothesis, so the function is likely piecewise. I need one more pair, preferably with neither input equal to {-}99, to identify the rule.","hypothesis":"unknown","confidence\_score":12\}} \\
            }
        }
    }
    & & \\ \\

    \theutterance \stepcounter{utterance}  
    & & & \multicolumn{4}{p{0.6\linewidth}}{
        \cellcolor[rgb]{0.9,0.9,0.9}{
            \makecell[{{p{\linewidth}}}]{
                \texttt{\tiny{[P1$\langle$GM]}}
                \texttt{OUTPUT: (({-}1, {-}1000), {-}4003)} \\
            }
        }
    }
     \\ \\

    \theutterance \stepcounter{utterance}  
    & \multicolumn{4}{p{0.6\linewidth}}{
        \cellcolor[rgb]{0.561,0.737,0.561}{
            \makecell[{{p{\linewidth}}}]{
                \texttt{\tiny{[P1$\rangle$GM]}}
                \texttt{\{"current\_mode":"confirming","input\_rationale":"The examples suggest a piecewise linear rule: if either argument is {-}99, the other argument is scaled by 102 and {-}99 contributes 3; otherwise the result is 4 times the smaller input plus 3 times the larger input. This matches all provided cases.","hypothesis":"def f(a: int, b: int) {-}> int:\textbackslash{}n\textbackslash{}t\textbackslash{}"\textbackslash{}"\textbackslash{}"Computes a piecewise linear combination of the two integers.\textbackslash{}"\textbackslash{}"\textbackslash{}"\textbackslash{}n\textbackslash{}tif a == {-}99 or b == {-}99:\textbackslash{}n\textbackslash{}t\textbackslash{}tother = b if a == {-}99 else a\textbackslash{}n\textbackslash{}t\textbackslash{}treturn 102 * other {-} 297\textbackslash{}n\textbackslash{}treturn 4 * min(a, b) + 3 * max(a, b)","confidence\_score":95\}} \\
            }
        }
    }
    & & \\ \\

    \theutterance \stepcounter{utterance}  
    & & & \multicolumn{4}{p{0.6\linewidth}}{
        \cellcolor[rgb]{0.9,0.9,0.9}{
            \makecell[{{p{\linewidth}}}]{
                \texttt{\tiny{[P1$\langle$GM]}}
                \texttt{OUTPUT: ((1, 99), 201)} \\
            }
        }
    }
     \\ \\

    \theutterance \stepcounter{utterance}  
    & \multicolumn{4}{p{0.6\linewidth}}{
        \cellcolor[rgb]{0.561,0.737,0.561}{
            \makecell[{{p{\linewidth}}}]{
                \texttt{\tiny{[P1$\rangle$GM]}}
                \texttt{\{"current\_mode":"probing","input\_rationale":"The examples with {-}99 reveal one special{-}case branch, but the general rule is still ambiguous. The new positive{-}positive example suggests the non{-}special branch may depend on the signs of the inputs, so I need one more mixed{-}sign example to pin it down.","hypothesis":"unknown","confidence\_score":18\}} \\
            }
        }
    }
    & & \\ \\

    \theutterance \stepcounter{utterance}  
    & & & \multicolumn{4}{p{0.6\linewidth}}{
        \cellcolor[rgb]{0.9,0.9,0.9}{
            \makecell[{{p{\linewidth}}}]{
                \texttt{\tiny{[P1$\langle$GM]}}
                \texttt{OUTPUT: (({-}1000, 2), {-}994)} \\
            }
        }
    }
     \\ \\

    \theutterance \stepcounter{utterance}  
    & \multicolumn{4}{p{0.6\linewidth}}{
        \cellcolor[rgb]{0.561,0.737,0.561}{
            \makecell[{{p{\linewidth}}}]{
                \texttt{\tiny{[P1$\rangle$GM]}}
                \texttt{\{"current\_mode":"probing","input\_rationale":"I need one more simple non{-}special case to determine whether the function depends on ordering, magnitude, or sign. A small positive{-}positive pair like (2,1) should distinguish the remaining candidates.","hypothesis":"unknown","confidence\_score":20\}} \\
            }
        }
    }
    & & \\ \\

    \theutterance \stepcounter{utterance}  
    & & & \multicolumn{4}{p{0.6\linewidth}}{
        \cellcolor[rgb]{0.9,0.9,0.9}{
            \makecell[{{p{\linewidth}}}]{
                \texttt{\tiny{[P1$\langle$GM]}}
                \texttt{OUTPUT: ((2, {-}99), {-}93)} \\
            }
        }
    }
     \\ \\

    \theutterance \stepcounter{utterance}  
    & \multicolumn{4}{p{0.6\linewidth}}{
        \cellcolor[rgb]{0.561,0.737,0.561}{
            \makecell[{{p{\linewidth}}}]{
                \texttt{\tiny{[P1$\rangle$GM]}}
                \texttt{\{"current\_mode":"confirming","input\_rationale":"The examples are consistent with a piecewise rule: a special case for any input equal to {-}99, otherwise coefficients depend on the sign pattern of the pair. This matches all seen outputs.","hypothesis":"def f(a: int, b: int) {-}> int:\textbackslash{}n\textbackslash{}t\textbackslash{}"\textbackslash{}"\textbackslash{}"Computes a piecewise linear combination of the two integers.\textbackslash{}"\textbackslash{}"\textbackslash{}"\textbackslash{}n\textbackslash{}tif a == {-}99 or b == {-}99:\textbackslash{}n\textbackslash{}t\textbackslash{}tother = b if a == {-}99 else a\textbackslash{}n\textbackslash{}t\textbackslash{}treturn 102 * other {-} 297\textbackslash{}n\textbackslash{}n\textbackslash{}tlo = a if a <= b else b\textbackslash{}n\textbackslash{}thi = b if a <= b else a\textbackslash{}n\textbackslash{}n\textbackslash{}tif a < 0 and b < 0:\textbackslash{}n\textbackslash{}t\textbackslash{}treturn 4 * lo + 3 * hi\textbackslash{}n\textbackslash{}tif a >= 0 and b >= 0:\textbackslash{}n\textbackslash{}t\textbackslash{}treturn 3 * lo + 2 * hi\textbackslash{}n\textbackslash{}treturn lo + 3 * hi","confidence\_score":84\}} \\
            }
        }
    }
    & & \\ \\

    \theutterance \stepcounter{utterance}  
    & & & \multicolumn{4}{p{0.6\linewidth}}{
        \cellcolor[rgb]{0.9,0.9,0.9}{
            \makecell[{{p{\linewidth}}}]{
                \texttt{\tiny{[P1$\langle$GM]}}
                \texttt{OUTPUT: (({-}99, {-}99), {-}10395)} \\
            }
        }
    }
     \\ \\

    \theutterance \stepcounter{utterance}  
    & \multicolumn{4}{p{0.6\linewidth}}{
        \cellcolor[rgb]{0.561,0.737,0.561}{
            \makecell[{{p{\linewidth}}}]{
                \texttt{\tiny{[P1$\rangle$GM]}}
                \texttt{```python} \\
\texttt{def f(a: int, b: int) {-}> int:} \\
\texttt{    """Computes a piecewise linear combination of the two integers."""} \\
\texttt{    if a == {-}99 or b == {-}99:} \\
\texttt{        other = b if a == {-}99 else a} \\
\texttt{        return 102 * other {-} 297} \\
\\ 
\texttt{    lo = a if a <= b else b} \\
\texttt{    hi = b if a <= b else a} \\
\\ 
\texttt{    if hi < 0:} \\
\texttt{        return 4 * lo + 3 * hi} \\
\texttt{    if lo >= 0:} \\
\texttt{        return 3 * lo + 2 * hi} \\
\texttt{    return lo + 3 * hi} \\
\texttt{```} \\
            }
        }
    }
    & & \\ \\

    \theutterance \stepcounter{utterance}  
    & & & \multicolumn{4}{p{0.6\linewidth}}{
        \cellcolor[rgb]{0.9,0.9,0.9}{
            \makecell[{{p{\linewidth}}}]{
                \texttt{\tiny{[P1$\langle$GM]}}
                \texttt{Game Violation: ParseError: Response is not valid JSON.} \\
\texttt{Output ONLY:} \\
\texttt{One JSON with confidence\_score, current\_mode, input\_rationale, and hypothesis keys ...} \\
\texttt{OR} \\
\texttt{SOLVE: ```python ... ```} \\
            }
        }
    }
     \\ \\

    \theutterance \stepcounter{utterance}  
    & \multicolumn{4}{p{0.6\linewidth}}{
        \cellcolor[rgb]{0.561,0.737,0.561}{
            \makecell[{{p{\linewidth}}}]{
                \texttt{\tiny{[P1$\rangle$GM]}}
                \texttt{SOLVE: ```python} \\
\texttt{def f(a: int, b: int) {-}> int:} \\
\texttt{    """Returns a piecewise linear combination of the two integers."""} \\
\texttt{    if a == {-}99 or b == {-}99:} \\
\texttt{        return 102 * (a + b) + 9801} \\
\\ 
\texttt{    lo = a if a <= b else b} \\
\texttt{    hi = b if a <= b else a} \\
\\ 
\texttt{    if hi < 0:} \\
\texttt{        return 4 * lo + 3 * hi} \\
\texttt{    if lo >= 0:} \\
\texttt{        return 3 * lo + 2 * hi} \\
\texttt{    return lo + 3 * hi} \\
\texttt{```} \\
            }
        }
    }
    & & \\ \\

    \theutterance \stepcounter{utterance}  
    & & & \multicolumn{2}{p{0.3\linewidth}}{
        \cellcolor[rgb]{0.9,0.9,0.9}{
            \makecell[{{p{\linewidth}}}]{
                \texttt{\tiny{[GM$|$GM]}}
                \texttt{Game Verdict: LOSS} \\
            }
        }
    }
    & & \\ \\

    \theutterance \stepcounter{utterance}  
    & & & \multicolumn{2}{p{0.3\linewidth}}{
        \cellcolor[rgb]{0.9,0.9,0.9}{
            \makecell[{{p{\linewidth}}}]{
                \texttt{\tiny{[GM$|$GM]}}
                \texttt{Incorrect.} \\
\texttt{Failed Cases:} \\
\texttt{Input: {[}{-}1, 99{]} | Expected: 393 | Got: 296} \\
\texttt{Input: {[}2, 0{]} | Expected: 6 | Got: 4} \\
\texttt{Input: {[}1, {-}10{]} | Expected: {-}17 | Got: {-}7} \\
            }
        }
    }
    & & \\ \\

    \theutterance \stepcounter{utterance}  
    & & & \multicolumn{2}{p{0.3\linewidth}}{
        \cellcolor[rgb]{0.9,0.9,0.9}{
            \makecell[{{p{\linewidth}}}]{
                \texttt{\tiny{[GM$|$GM]}}
                \texttt{def triple\_sum\_minus\_product(a: int, b: int) {-}> int:} \\
\texttt{    return 3 * (a + b) {-} a * b} \\
            }
        }
    }
    & & \\ \\

    \theutterance \stepcounter{utterance}  
    & & & \multicolumn{2}{p{0.3\linewidth}}{
        \cellcolor[rgb]{0.9,0.9,0.9}{
            \makecell[{{p{\linewidth}}}]{
                \texttt{\tiny{[GM$|$GM]}}
                \texttt{def f(a: int, b: int) {-}> int:} \\
\texttt{    """Returns a piecewise linear combination of the two integers."""} \\
\texttt{    if a == {-}99 or b == {-}99:} \\
\texttt{        return 102 * (a + b) + 9801} \\
\\ 
\texttt{    lo = a if a <= b else b} \\
\texttt{    hi = b if a <= b else a} \\
\\ 
\texttt{    if hi < 0:} \\
\texttt{        return 4 * lo + 3 * hi} \\
\texttt{    if lo >= 0:} \\
\texttt{        return 3 * lo + 2 * hi} \\
\texttt{    return lo + 3 * hi} \\
            }
        }
    }
    & & \\ \\

    \theutterance \stepcounter{utterance}  
    & & & \multicolumn{2}{p{0.3\linewidth}}{
        \cellcolor[rgb]{0.9,0.9,0.9}{
            \makecell[{{p{\linewidth}}}]{
                \texttt{\tiny{[GM$|$GM]}}
                \texttt{Mode: passive\_examples} \\
\texttt{Accuracy: 0.360} \\
\texttt{Information requests: 9/15} \\
\texttt{Revealed passive examples: 10} \\
\texttt{Efficiency: 0.467} \\
\texttt{Slack: 6} \\
\texttt{Observed{-}pair consistency: 1.000} \\
\texttt{Parse errors: 1} \\
\texttt{Runtime errors: 0} \\
            }
        }
    }
    & & \\ \\

\end{supertabular}
}

\end{document}
